\editionPage[\pdfbookmark[0]{Corrigés des défis}{bm-corriges-defis}]{defi-cor-1-1}{D01 · Corrigé}{Le dernier souffle de César}{\hyperref[defi-ex-1]{Énoncé p.~\pageref*{defi-ex-1}} · Démarche et vérifications}
\begin{blocEdition}{defi-cor-1-1}
\editionEtape{1. L'atmosphère en moles}
La masse de l'atmosphère vaut $\qty{5,1e18}{kg}=\qty{5,1e21}{g}$, d'où
\[n_\text{atm}=\frac{m}{M}=\frac{\num{5.1e21}}{29}=\res{\qty{1,76e20}{mol}}\qquad\text{et}\qquad N_\text{atm}=n_\text{atm}\times N_\mathrm{A}=\res{\num{1.06e44}}\ \text{molécules}.\]
\editionEtape{2. Une inspiration en moles}
\[n_\text{souffle}=\frac{V}{V_\mathrm{m}}=\frac{0{,}50}{24{,}0}=\res{\qty{2,08e-2}{mol}},\]
d'où $N_\text{souffle}=n_\text{souffle}\times N_\mathrm{A}=\num{2.08e-2}\times\num{6.02e23}=\res{\num{1.25e22}}$ molécules.
\editionEtape{3. Deux nombres gigantesques et voisins}
L'atmosphère contient $\dfrac{n_\text{atm}}{n_\text{souffle}}=\dfrac{\num{1.76e20}}{\num{2.08e-2}}=\res{\num{8.4e21}}$ inspirations, soit un peu \textbf{moins} que les
$\num{1.25e22}$ molécules d'une seule inspiration. Ce sont deux nombres considérables, du même ordre de grandeur : c'est
la clé du problème.
\editionEtape{4. Combien de molécules de César par inspiration ?}
Le dernier souffle de César représente la fraction $f=\dfrac{1}{\num{8.4e21}}=\num{1.18e-22}$ de l'atmosphère. Le mélange étant supposé
uniforme, chacune de nos inspirations contient la même fraction de ses molécules :
\[N_\text{César}=f\times N_\text{souffle}=\num{1.18e-22}\times\num{1.25e22}\approx\res{1{,}5}.\]
En moyenne, \textbf{chaque inspiration contient une ou deux molécules} du dernier souffle de César, comme de n'importe
quelle expiration passée.
\end{blocEdition}

\editionPage{defi-cor-1-2}{D01 · Corrigé (suite)}{Le dernier souffle de César}{\hyperref[defi-ex-1]{Énoncé p.~\pageref*{defi-ex-1}} · Démarche et vérifications}
\begin{blocEdition}{defi-cor-1-2}
\label{figedition-defi-cor-1-2-1}\schema[corrige]{defi-01-cor}
\editionEtape{5. En une heure}
Quinze inspirations par minute pendant soixante minutes font 900 inspirations : on respire environ
$1{,}5\times900\approx\res{\num{1.3e3}}$ molécules du dernier souffle de César en une heure.
\editionEtape{6. Discuter les hypothèses}
Le dioxygène inspiré par César a pu être consommé, puis renouvelé ; une partie du dioxyde de carbone expiré s'est
dissoute dans les océans. Mais le diazote, qui forme près des quatre cinquièmes de l'air, réagit très peu, et
deux mille ans suffisent largement à mélanger toute l'atmosphère, ce qui prend seulement quelques années. Même si
seules les molécules de diazote comptaient, il en resterait environ $0{,}78\times1{,}5\approx1{,}2$ par inspiration : l'ordre de
grandeur, \res{\text{une molécule par inspiration}}, ne change pas.
\editionRetenir{La constante d'Avogadro est si grande qu'une petite quantité de matière contient un nombre immense d'entités :
une seule inspiration renferme plus de molécules que l'atmosphère entière ne contient d'inspirations. Avant de
conclure à l'impossible, il faut calculer.}
\end{blocEdition}

\editionPage{defi-cor-2-1}{D02 · Corrigé}{Deux colorants, une seule couleur}{\hyperref[defi-ex-2]{Énoncé p.~\pageref*{defi-ex-2}} · Démarche et vérifications}
\begin{blocEdition}{defi-cor-2-1}
\editionEtape{1. Une mesure, deux inconnues}
À \qty{427}{nm}, les deux colorants absorbent : l'absorbance mesurée est la somme de deux termes, et une seule équation ne
permet pas de déterminer deux concentrations. Il faut une \res{\text{seconde mesure}}, à une autre longueur d'onde.
\editionEtape{2. Additivité des absorbances}
\[A_{427}=0{,}0270\,c_\mathrm{J}+0{,}0060\,c_\mathrm{B}\qquad\qquad A_{638}=0\times c_\mathrm{J}+0{,}106\,c_\mathrm{B}=0{,}106\,c_\mathrm{B},\]
avec les concentrations en \unit{\micro mol.L^{-1}}.
\editionEtape{3. Résolution}
La seconde équation donne directement $c_\mathrm{B}=\dfrac{0{,}53}{0{,}106}=\res{\qty{5,0}{\micro mol.L^{-1}}}$. En reportant dans la première :
\[c_\mathrm{J}=\frac{A_{427}-0{,}0060\,c_\mathrm{B}}{0{,}0270}=\frac{0{,}62-0{,}030}{0{,}0270}=\res{\qty{21,9}{\micro mol.L^{-1}}}.\]
\editionEtape{4. L'erreur commise en négligeant le bleu}
Sans tenir compte du bleu patenté, on trouverait $\dfrac{0{,}62}{0{,}0270}=\qty{23,0}{\micro mol.L^{-1}}$, soit une erreur relative de
$\dfrac{23{,}0-21{,}9}{21{,}9}\approx\res{\qty{5}{\percent}}$ : le bleu patenté ne contribue que pour 0,030 à l'absorbance de 0,62, mais l'erreur
n'est pas négligeable pour un contrôle.
\end{blocEdition}

\editionPage{defi-cor-2-2}{D02 · Corrigé (suite)}{Deux colorants, une seule couleur}{\hyperref[defi-ex-2]{Énoncé p.~\pageref*{defi-ex-2}} · Démarche et vérifications}
\begin{blocEdition}{defi-cor-2-2}
\label{figedition-defi-cor-2-2-1}\schema[corrige]{02-ex11}
\editionEtape{5. Dans le sirop}
Le sirop est vingt fois plus concentré : $c_\mathrm{J}=20\times21{,}9\approx\res{\qty{438}{\micro mol.L^{-1}}}$ et $c_\mathrm{B}=20\times5{,}0=\res{\qty{100}{\micro mol.L^{-1}}}$.
\editionEtape{6. Le cas général}
Si les deux colorants absorbent aux deux longueurs d'onde $\lambda_1$ et $\lambda_2$ :
\[A_1=k_{\mathrm{J}1}\,c_\mathrm{J}+k_{\mathrm{B}1}\,c_\mathrm{B}\qquad\qquad A_2=k_{\mathrm{J}2}\,c_\mathrm{J}+k_{\mathrm{B}2}\,c_\mathrm{B}.\]
On exprime $c_\mathrm{B}$ à partir de la seconde équation, on le reporte dans la première, qui ne contient plus que $c_\mathrm{J}$, puis on
revient à $c_\mathrm{B}$. Si les deux colorants absorbaient de la même façon aux deux longueurs d'onde, les deux équations
seraient presque proportionnelles : la moindre erreur de mesure fausserait fortement le résultat. On choisit donc les
maximums d'absorption de chaque colorant, où leurs contributions sont \res{\text{les plus différentes}}.
\editionRetenir{Les absorbances de plusieurs espèces s'ajoutent. Avec autant de longueurs d'onde que d'espèces absorbantes,
on obtient un système d'équations du premier degré qui donne toutes les concentrations.}
\end{blocEdition}

\editionPage{defi-cor-3-1}{D03 · Corrigé}{L'arbre de Diane prend du poids}{\hyperref[defi-ex-3]{Énoncé p.~\pageref*{defi-ex-3}} · Démarche et vérifications}
\begin{blocEdition}{defi-cor-3-1}
\editionEtape{1. L'équation}
Le cuivre est oxydé, $\ce{Cu -> Cu^{2+} + 2e-}$, et les ions argent sont réduits, $\ce{Ag+ + e- -> Ag}$, à multiplier par 2 :
\[\res{\ce{Cu + 2Ag+ -> Cu^{2+} + 2Ag}}.\]
\editionEtape{2. Le réactif épuisé}
$n(\ce{Ag+})=0{,}100\times0{,}100=\qty{0,0100}{mol}$ et $n(\ce{Cu})=\dfrac{5{,}00}{63{,}5}=\qty{0,0787}{mol}$. Les ions argent ne peuvent
oxyder que $\dfrac{0{,}0100}{2}=\qty{5,00e-3}{mol}$ de cuivre, bien moins que la quantité présente : les \res{\text{ions argent}} seront
entièrement consommés, le cuivre est en excès.
\editionEtape{3. La masse finale du fil}
Masse de cuivre consommée : $\num{5.00e-3}\times63{,}5=\res{\qty{0,318}{g}}$. Masse d'argent déposée : $0{,}0100\times107{,}9=\res{\qty{1,08}{g}}$.
Le fil gagne $1{,}079-0{,}318\approx\qty{0,76}{g}$ et pèse finalement $\res{\qty{5,76}{g}}$. Il s'alourdit alors qu'il est rongé :
l'argent est plus lourd que le cuivre, et deux atomes d'argent remplacent chaque atome de cuivre.
\editionEtape{4. La masse en fonction de l'avancement}
Lorsque $x$ mol de cuivre ont réagi, le fil a perdu $63{,}5\,x$ grammes de cuivre et gagné $2x$ moles d'argent, soit
$2\times107{,}9\,x=215{,}8\,x$ grammes. Ainsi
\[m=5{,}00-63{,}5\,x+215{,}8\,x=\res{5{,}00+152{,}3\,x}.\]
Le nombre $152{,}3=2\times107{,}9-63{,}5$ est la \res{\text{prise de masse du fil par mole de cuivre ayant réagi}}, en \unit{g.mol^{-1}}. La
masse du fil augmente proportionnellement à $x$, jusqu'à $x=\qty{5,00e-3}{mol}$.
\end{blocEdition}

\editionPage{defi-cor-3-2}{D03 · Corrigé (suite)}{L'arbre de Diane prend du poids}{\hyperref[defi-ex-3]{Énoncé p.~\pageref*{defi-ex-3}} · Démarche et vérifications}
\begin{blocEdition}{defi-cor-3-2}
\label{figedition-defi-cor-3-2-1}\schema[corrige]{defi-03-cor}
\editionEtape{5. Une pesée en cours de réaction}
La prise de masse vaut $5{,}38-5{,}00=\qty{0,38}{g}$, donc $x=\dfrac{0{,}38}{152{,}3}=\res{\qty{2,5e-3}{mol}}$ : la moitié des ions argent a réagi.
Il s'est formé autant d'ions cuivre(II) que de cuivre consommé : $[\ce{Cu^{2+}}]=\dfrac{\num{2.5e-3}}{0{,}100}=\res{\qty{0,025}{mol.L^{-1}}}$. En fin
de réaction, elle atteindra \qty{0,0500}{mol.L^{-1}}.
\editionEtape{6. Dissoudre tout le fil}
Il faudrait $2\times0{,}0787=\qty{0,157}{mol}$ d'ions argent, soit un volume $V=\dfrac{0{,}157}{0{,}100}\approx\res{\qty{1,57}{L}}$ de solution,
plus de quinze fois le volume utilisé.
\editionEtape{7. Suivre la réaction sans peser}
Les ions cuivre(II) sont les seules espèces colorées : on peut mesurer régulièrement l'\res{\text{absorbance}} de la solution
vers \qty{800}{nm}, où elle est maximale, et en déduire $[\ce{Cu^{2+}}]$ grâce à une droite d'étalonnage, puis $x$.
\editionRetenir{Les quantités qui réagissent sont dans les proportions des nombres stœchiométriques. Une grandeur mesurable,
masse ou absorbance, qui varie proportionnellement à l'avancement permet de suivre une transformation.}
\end{blocEdition}

\editionPage{defi-cor-4-1}{D04 · Corrigé}{La sonde lambda}{\hyperref[defi-ex-4]{Énoncé p.~\pageref*{defi-ex-4}} · Démarche et vérifications}
\begin{blocEdition}{defi-cor-4-1}
\editionEtape{1. L'équation de combustion}
Huit atomes de carbone donnent 8 molécules de $\ce{CO2}$, dix-huit atomes d'hydrogène 9 molécules d'eau ; il faut alors
$\dfrac{16+9}{2}=12{,}5$ molécules de dioxygène. En multipliant par 2 :
\[\res{\ce{2C8H18 + 25O2 -> 16CO2 + 18H2O}}.\]
\editionEtape{2. L'air d'un mélange stœchiométrique}
Une mole d'octane réagit avec $\dfrac{25}{2}=\res{\qty{12,5}{mol}}$ de dioxygène. L'air en contient 21~\%, d'où
\[n_\text{air}=\frac{12{,}5}{0{,}21}=\qty{59,5}{mol}\qquad\text{et}\qquad m_\text{air}=59{,}5\times29{,}0\approx\res{\qty{1726}{g}}.\]
\editionEtape{3. Le rapport air/essence}
$M(\ce{C8H18})=8\times12{,}0+18\times1{,}0=\qty{114,0}{g.mol^{-1}}$, d'où $\dfrac{1726}{114{,}0}\approx\res{15{,}1}$ : il faut environ \qty{15}{kg} d'air pour brûler
\qty{1}{kg} d'octane. L'essence réelle est un mélange qui contient aussi des molécules plus pauvres en hydrogène,
comme le toluène $\ce{C7H8}$ : elles consomment moins de dioxygène par kilogramme, d'où un rapport \res{\text{un peu plus faible}}.
\editionEtape{4. Mélange riche, mélange pauvre}
Pour $\lambda=0{,}90$ : $n(\ce{O2})=0{,}90\times12{,}5=\qty{11,25}{mol}$. Les quotients valent $\dfrac{1{,}0}{2}=0{,}50$ et $\dfrac{11{,}25}{25}=0{,}45$ : le
\res{\text{dioxygène est limitant}}, $x_\mathrm{max}=\qty{0,45}{mol}$. Le programme, avec \code{reactifs = \{'C8H18': [2, 1.0], 'O2': [25, 11.25]\}},
affiche \qty{0,1}{mol} d'octane restant, \qty{7,2}{mol} de $\ce{CO2}$ et \qty{8,1}{mol} d'eau.\par\smallskip
Pour $\lambda=1{,}10$ : $n(\ce{O2})=\qty{13,75}{mol}$, $\dfrac{13{,}75}{25}=0{,}55>0{,}50$ : l'\res{\text{octane est limitant}}, $x_\mathrm{max}=\qty{0,50}{mol}$ ; il reste
$13{,}75-25\times0{,}50=\qty{1,25}{mol}$ de dioxygène, et il se forme \qty{8,0}{mol} de $\ce{CO2}$ et \qty{9,0}{mol} d'eau.
\end{blocEdition}

\editionPage{defi-cor-4-2}{D04 · Corrigé (suite)}{La sonde lambda}{\hyperref[defi-ex-4]{Énoncé p.~\pageref*{defi-ex-4}} · Démarche et vérifications}
\begin{blocEdition}{defi-cor-4-2}
\label{figedition-defi-cor-4-2-1}\schema[corrige]{defi-04-cor}
\editionEtape{5. Les imbrûlés}
Dans un mélange riche, le dioxygène est le réactif limitant : une partie du carburant, ici \qty{10}{\percent} pour $\lambda=0{,}90$, ne
peut pas brûler complètement. Elle sort du moteur sous forme d'\res{\text{imbrûlés}}, ou ne s'oxyde qu'en partie en
\res{\text{monoxyde de carbone}}, faute de dioxygène. Dans un mélange pauvre, c'est au contraire le dioxygène en excès qui
empêche le pot catalytique d'éliminer les oxydes d'azote : d'où la régulation autour de $\lambda=1$.
\editionEtape{6. Le dioxyde de carbone émis}
Par kilomètre, la voiture brûle $0{,}060\times750=\qty{45,0}{g}$ d'essence, soit $\dfrac{45{,}0}{114{,}0}=\qty{0,395}{mol}$ d'octane. Une mole d'octane
produit 8 moles de $\ce{CO2}$ :
\[m(\ce{CO2})=8\times0{,}395\times44{,}0\approx\res{\qty{139}{g}}\text{ par kilomètre}.\]
C'est l'ordre de grandeur affiché pour une voiture de cette consommation : chaque litre d'essence brûlé rejette environ
\qty{2,3}{kg} de dioxyde de carbone.
\editionRetenir{Le rapport entre les quantités de réactifs introduites et les proportions stœchiométriques décide du réactif
limitant, donc de ce qui reste dans l'état final : imbrûlés pour un mélange riche, dioxygène pour un mélange pauvre.}
\end{blocEdition}

\editionPage{defi-cor-5-1}{D05 · Corrigé}{Que vaut vraiment une eau de Javel ?}{\hyperref[defi-ex-5]{Énoncé p.~\pageref*{defi-ex-5}} · Démarche et vérifications}
\begin{blocEdition}{defi-cor-5-1}
\editionEtape{1. La réaction des ions hypochlorite}
De chaque côté : 1 Cl, 1 O, 2 I, 2 H et une charge $-1$. L'ion hypochlorite est l'\res{\text{oxydant}}, réduit en ion chlorure
($\ce{ClO- + 2H+ + 2e- -> Cl- + H2O}$) ; l'ion iodure est le \res{\text{réducteur}}, oxydé en diiode ($\ce{2I- -> I2 + 2e-}$).
\editionEtape{2. Un réactif en excès}
Tous les ions hypochlorite doivent être transformés en diiode : les ions iodure ne doivent pas être limitants. Pour le
flacon neuf, il faut $2\times\num{3.65e-4}=\qty{7,30e-4}{mol}$ d'ions iodure ; \qty{1}{g} de $\ce{KI}$ en apporte $\dfrac{1}{166{,}0}\approx\qty{6,0e-3}{mol}$,
soit plus de huit fois plus : \res{\text{l'excès est largement assuré}}.
\editionEtape{3. La chaîne des quantités}
À l'équivalence du titrage, $n(\ce{I2})=\dfrac{c\,V_\mathrm{E}}{2}$ ; la première réaction donne $n(\ce{I2})=n(\ce{ClO-})$. Ainsi
\[\res{n(\ce{ClO-})=\frac{c\,V_\mathrm{E}}{2}}=\frac{\num{5.00e-2}\times\num{14.6e-3}}{2}=\qty{3,65e-4}{mol}\text{ dans \qty{10,0}{mL} de solution diluée.}\]
\editionEtape{4. La concentration de l'eau de Javel}
Solution diluée : $\dfrac{\num{3.65e-4}}{\num{10.0e-3}}=\qty{3,65e-2}{mol.L^{-1}}$ ; eau de Javel, dix fois plus concentrée : $\res{\qty{0,365}{mol.L^{-1}}}$.
\editionEtape{5. Le chlore actif}
Un litre libérerait \qty{0,365}{mol} de dichlore, soit $0{,}365\times71{,}0\approx\qty{25,9}{g}$ : \qty{2,59}{g} pour \qty{100}{mL}, soit $\res{\qty{2,6}{\percent}}$.
L'étiquette est \res{\text{exacte}}.
\end{blocEdition}

\editionPage{defi-cor-5-2}{D05 · Corrigé (suite)}{Que vaut vraiment une eau de Javel ?}{\hyperref[defi-ex-5]{Énoncé p.~\pageref*{defi-ex-5}} · Démarche et vérifications}
\begin{blocEdition}{defi-cor-5-2}
\label{figedition-defi-cor-5-2-1}\schema[corrige]{defi-05-cor}
\editionEtape{6. Le vieux flacon}
La concentration est proportionnelle au volume équivalent : $\dfrac{0{,}365\times11{,}2}{14{,}6}\approx\qty{0,280}{mol.L^{-1}}$, soit \qty{19,9}{g} de chlore
actif par litre, environ \qty{2,0}{\percent}. Le flacon a perdu $1-\dfrac{11{,}2}{14{,}6}\approx\res{\qty{23}{\percent}}$ de son pouvoir oxydant. Les ions
hypochlorite se décomposent lentement, d'autant plus vite que la température est élevée et que le flacon est exposé à la
lumière : il faut \res{\text{conserver l'eau de Javel au frais et à l'abri de la lumière}}, et l'utiliser dans les mois qui suivent
sa fabrication.
\editionRetenir{Dans un titrage indirect, l'espèce à doser réagit d'abord avec un réactif en excès ; on titre ensuite le
produit formé, puis on remonte la chaîne des quantités de matière en respectant les proportions de chaque équation.}
\end{blocEdition}

\editionPage{defi-cor-6-1}{D06 · Corrigé}{La molécule d'eau en géométrie}{\hyperref[defi-ex-6]{Énoncé p.~\pageref*{defi-ex-6}} · Démarche et vérifications}
\begin{blocEdition}{defi-cor-6-1}
\editionEtape{1. Deux distances dans le cube}
La diagonale d'une face carrée d'arête 2 vaut $\sqrt{2^2+2^2}=\res{2\sqrt2}$ : c'est la distance H--H. La grande diagonale du cube
vaut $\sqrt{\left(2\sqrt2\right)^2+2^2}=\sqrt{12}=2\sqrt3$ ; le carbone est en son milieu, donc C--H $=\res{\sqrt3}$.
\editionEtape{2. L'angle tétraédrique}
Dans le triangle formé par le carbone et deux atomes d'hydrogène :
\[\left(2\sqrt2\right)^2=\left(\sqrt3\right)^2+\left(\sqrt3\right)^2-2\times3\cos\theta\quad\Longleftrightarrow\quad8=6-6\cos\theta\quad\Longleftrightarrow\quad\res{\cos\theta=-\frac13},\]
d'où $\res{\theta\approx\qty{109,5}{\degree}}$ : c'est l'angle entre deux sommets d'un tétraèdre régulier vus depuis son centre.
\editionEtape{3. L'effet des doublets non liants}
Dans l'ammoniac et dans l'eau, les doublets s'orientent encore vers les sommets d'un tétraèdre, mais les doublets non
liants, qui n'appartiennent qu'à l'atome central, occupent plus de place et \res{\text{repoussent davantage}} les doublets liants.
Plus il y a de doublets non liants, plus l'angle diminue : \qty{109,5}{\degree}, puis \qty{107}{\degree}, puis \qty{104,5}{\degree}.
\editionEtape{4. La somme de deux vecteurs}
Les deux vecteurs ont la même norme $p_1$ et font chacun un angle $\dfrac{\alpha}{2}$ avec la bissectrice. Leurs composantes
perpendiculaires à la bissectrice s'annulent ; leurs composantes le long de la bissectrice valent chacune $p_1\cos\left(\dfrac{\alpha}{2}\right)$.
Ainsi $\res{p=2\,p_1\cos\left(\dfrac{\alpha}{2}\right)}$.
\end{blocEdition}

\editionPage{defi-cor-6-2}{D06 · Corrigé (suite)}{La molécule d'eau en géométrie}{\hyperref[defi-ex-6]{Énoncé p.~\pageref*{defi-ex-6}} · Démarche et vérifications}
\begin{blocEdition}{defi-cor-6-2}
\label{figedition-defi-cor-6-2-1}\schema[corrige]{defi-06-cor}
\editionEtape{5. La molécule d'eau}
$p=2\times1{,}51\times\cos(\qty{52,25}{\degree})\approx2\times1{,}51\times0{,}612\approx\res{\qty{1,85}{D}}$ : la valeur calculée coïncide avec la
mesure. Le modèle des liaisons polarisées, associé à la géométrie coudée, rend compte de la polarité de l'eau.
\editionEtape{6. Le rôle de la géométrie}
Pour une molécule linéaire, $\alpha=\qty{180}{\degree}$ et $\cos(\qty{90}{\degree})=0$ : $\res{p=0}$, la molécule serait apolaire malgré ses liaisons
polarisées. Pour $\alpha=\qty{109,5}{\degree}$, on obtiendrait environ \qty{1,74}{D}. Dans le dioxyde de carbone, linéaire, les deux vecteurs
C--O sont opposés : $\res{p=0}$. Une molécule est polaire si la somme des vecteurs de ses liaisons n'est pas nulle :
c'est la traduction vectorielle du critère « G$^{+}$ et G$^{-}$ distincts ».
\editionRetenir{La polarité d'une molécule dépend de la polarisation de ses liaisons et de leur disposition : deux liaisons
identiques se compensent exactement si elles sont opposées, partiellement si elles forment un angle.}
\end{blocEdition}

\editionPage{defi-cor-7-1}{D07 · Corrigé}{L'étiquette d'une eau minérale}{\hyperref[defi-ex-7]{Énoncé p.~\pageref*{defi-ex-7}} · Démarche et vérifications}
\begin{blocEdition}{defi-cor-7-1}
\editionEtape{1. Les concentrations}
$M(\ce{SO4^{2-}})=\qty{96,1}{g.mol^{-1}}$, $M(\ce{HCO3-})=\qty{61,0}{g.mol^{-1}}$, $M(\ce{NO3-})=\qty{62,0}{g.mol^{-1}}$. En divisant chaque concentration en masse par la masse
molaire, on obtient, en \unit{mmol.L^{-1}} :
\begin{center}
\begin{tabular}{lcccc@{\qquad}lcccc}
\toprule
 & $\ce{Ca^{2+}}$ & $\ce{Mg^{2+}}$ & $\ce{Na+}$ & $\ce{K+}$ & & $\ce{SO4^{2-}}$ & $\ce{HCO3-}$ & $\ce{Cl-}$ & $\ce{NO3-}$ \\
\midrule
$c$ & 12,1 & 3,46 & 0,617 & 0,105 & $c$ & 12,4 & 6,61 & 0,394 & 0,069 \\
$|z|\,c$ & 24,2 & 6,91 & 0,617 & 0,105 & $|z|\,c$ & 24,7 & 6,61 & 0,394 & 0,069 \\
\bottomrule
\end{tabular}
\end{center}
\editionEtape{2. Les deux quantités de charges}
Charges positives : $24{,}2+6{,}91+0{,}617+0{,}105\approx\res{\qty{31,9}{mmol.L^{-1}}}$. Charges négatives : $24{,}7+6{,}61+0{,}394+0{,}069\approx\res{\qty{31,8}{mmol.L^{-1}}}$.
\editionEtape{3. L'électroneutralité}
L'écart relatif vaut environ $\dfrac{31{,}9-31{,}8}{31{,}9}\approx\res{\qty{0,3}{\percent}}$ : compte tenu des arrondis de l'étiquette et des ions non dosés, en
très faible quantité, l'étiquette est \res{\text{cohérente}}.
\editionEtape{4. Une eau qui a traversé du gypse}
Les ions calcium (\qty{12,1}{mmol.L^{-1}}) et sulfate (\qty{12,4}{mmol.L^{-1}}) sont en quantités presque égales : leur rapport vaut 0,98. C'est la
signature de la dissolution du \res{\text{gypse}}, $\ce{CaSO4}$, $2\,\ce{H2O}$ $\ce{->}$ $\ce{Ca^{2+}(aq)}+\ce{SO4^{2-}(aq)}$ (+ eau), qui libère autant d'ions de
chaque sorte. Les ions magnésium et hydrogénocarbonate proviennent d'autres roches, comme la dolomie.
\editionEtape{5. La masse de gypse dissous}
Au moins \qty{12,1}{mmol} de gypse par litre, de masse molaire $40{,}1+96{,}1+2\times18{,}0=\qty{172,2}{g.mol^{-1}}$ : $m\approx\res{\qty{2,09}{g}}$ par litre.
\end{blocEdition}

\editionPage{defi-cor-7-2}{D07 · Corrigé (suite)}{L'étiquette d'une eau minérale}{\hyperref[defi-ex-7]{Énoncé p.~\pageref*{defi-ex-7}} · Démarche et vérifications}
\begin{blocEdition}{defi-cor-7-2}
\label{figedition-defi-cor-7-2-1}\schema[corrige]{defi-07-cor}
\editionEtape{6. Un programme de contrôle}
\begin{python}
cations = {'Ca2+': (486, 40.1, 2), 'Mg2+': (84, 24.3, 2),
           'Na+': (14.2, 23.0, 1), 'K+': (4.1, 39.1, 1)}
anions = {'SO42-': (1187, 96.1, 2), 'HCO3-': (403, 61.0, 1),
          'Cl-': (14, 35.5, 1), 'NO3-': (4.3, 62.0, 1)}

def charges(ions):
    total = 0
    for t, M, z in ions.values():      # t en mg/L, M en g/mol
        total = total + z * t / M      # mmol de charges par litre
    return total

plus, moins = charges(cations), charges(anions)
ecart = 100 * (plus - moins) / plus
print(round(plus, 1), round(moins, 1), round(ecart, 1))
\end{python}
\begin{console}
31.9 31.8 0.3
\end{console}
\editionRetenir{Une solution est électriquement neutre : la somme des $|z|\,c$ des cations est égale à celle des anions. Ce bilan permet
de contrôler une analyse ou une étiquette.}
\end{blocEdition}

\editionPage{defi-cor-8-1}{D08 · Corrigé}{L'odeur du chocolat}{\hyperref[defi-ex-8]{Énoncé p.~\pageref*{defi-ex-8}} · Démarche et vérifications}
\begin{blocEdition}{defi-cor-8-1}
\editionEtape{1. Aldéhyde ou cétone}
Un seul atome d'oxygène, engagé dans un groupe carbonyle : la molécule ne peut être ni un alcool (pas de $\ce{-OH}$) ni un acide
carboxylique (deux atomes d'oxygène). Si le carbone du groupe C=O est lié à un hydrogène, c'est un \res{\text{aldéhyde}} ; s'il est
lié à deux atomes de carbone, une \res{\text{cétone}}.
\editionEtape{2. Les sept candidats}
\label{figedition-defi-cor-8-1-1}\schema[corrige]{defi-08-cor}
\editionEtape{3. L'indice A}
Le spectre présente une bande forte et fine vers \qty{1728}{cm^{-1}} (C=O), des bandes C--H vers 2960 et \qty{2875}{cm^{-1}}, et surtout
\res{\text{deux bandes moyennes}} vers 2715 et \qty{2810}{cm^{-1}} : ce sont les bandes C--H d'un \res{\text{aldéhyde}}. Aucune bande O--H n'apparaît.
Il reste quatre candidats : les aldéhydes.
\end{blocEdition}

\editionPage{defi-cor-8-2}{D08 · Corrigé (suite)}{L'odeur du chocolat}{\hyperref[defi-ex-8]{Énoncé p.~\pageref*{defi-ex-8}} · Démarche et vérifications}
\begin{blocEdition}{defi-cor-8-2}
\editionEtape{4. Les indices B et C}
L'indice B écarte le pentanal, dont la chaîne est linéaire. Dans le 2-méthylbutanal et le 2,2-diméthylpropanal, la
ramification est portée par le carbone numéro 2, voisin du groupe $\ce{-CHO}$ : l'indice C les écarte. Il reste
\[\res{\ce{CH3-CH(CH3)-CH2-CHO}},\]
le \res{\text{3-méthylbutanal}} : chaîne principale de quatre carbones dont celui du groupe aldéhyde (numéro 1), groupe méthyle sur le
carbone 3.
\editionEtape{5. La masse molaire}
$M(\ce{C5H10O})=5\times12{,}0+10\times1{,}0+16{,}0=\res{\qty{86,0}{g.mol^{-1}}}$.
\editionEtape{6. Le rôle de chaque indice}
Chaque indice élimine une partie des candidats : le spectre en écarte trois, l'indice B un, l'indice C deux. Seule leur
combinaison laisse une molécule. Si l'on savait que la chaîne principale compte quatre carbones, les seuls aldéhydes
possibles seraient le 2-méthylbutanal et le 3-méthylbutanal : l'\res{\text{indice B}} deviendrait inutile, car ces deux molécules
sont ramifiées.
\editionRetenir{Pour identifier une molécule, on dresse la liste complète des candidats compatibles avec la formule brute,
puis on élimine à chaque indice ; le spectre infrarouge désigne la famille, pas la molécule.}
\end{blocEdition}

\editionPage{defi-cor-9-1}{D09 · Corrigé}{Deux procédés pour l'ibuprofène}{\hyperref[defi-ex-9]{Énoncé p.~\pageref*{defi-ex-9}} · Démarche et vérifications}
\begin{blocEdition}{defi-cor-9-1}
\editionEtape{1. Le rendement global}
Si l'on part de $n_0$, la première étape donne $\eta_1n_0$, la deuxième $\eta_2(\eta_1n_0)$, et ainsi de suite. Après $k$ étapes :
$n=\eta_1\eta_2\cdots\eta_k\,n_0$ : le rendement global est le \res{\text{produit des rendements}}.
\editionEtape{2. Six étapes contre trois}
Procédé en six étapes : $0{,}90^6\approx\res{\qty{53}{\percent}}$. Procédé en trois étapes : $0{,}90^3\approx\res{\qty{73}{\percent}}$. Raccourcir la synthèse fait gagner
vingt points de rendement.
\editionEtape{3. L'économie d'atomes}
$M(\ce{C13H18O2})=13\times12{,}0+18\times1{,}0+2\times16{,}0=\qty{206,0}{g.mol^{-1}}$. Les réactifs ont pour masses molaires 134,0 ; 102,0 ; 2,0 et
28,0, soit \qty{266,0}{g.mol^{-1}} au total. L'économie d'atomes vaut $\dfrac{206{,}0}{266{,}0}\approx\res{\qty{77}{\percent}}$, contre \qty{40}{\percent} pour l'ancien procédé.
\editionEtape{4. Les déchets incompressibles}
Ancien procédé : $\dfrac{1}{0{,}40}=\qty{2,5}{kg}$ de réactifs pour \qty{1}{kg} de produit, soit \res{\qty{1,5}{kg}} de déchets. Nouveau procédé :
$\dfrac{266{,}0}{206{,}0}\approx\qty{1,29}{kg}$ de réactifs, soit \res{\qty{0,29}{kg}} de déchets ; ce sous-produit est d'ailleurs l'acide éthanoïque,
récupéré et valorisé.
\editionEtape{5. Avec les rendements réels}
Ancien procédé : $\dfrac{2{,}5}{0{,}53}\approx\res{\qty{4,7}{kg}}$ de réactifs. Nouveau procédé : $\dfrac{1{,}29}{0{,}729}\approx\res{\qty{1,77}{kg}}$ de réactifs, près de trois fois
moins.
\end{blocEdition}

\editionPage{defi-cor-9-2}{D09 · Corrigé (suite)}{Deux procédés pour l'ibuprofène}{\hyperref[defi-ex-9]{Énoncé p.~\pageref*{defi-ex-9}} · Démarche et vérifications}
\begin{blocEdition}{defi-cor-9-2}
\label{figedition-defi-cor-9-2-1}\schema[corrige]{defi-09-cor}
\editionEtape{6. Un programme et une conclusion}
\begin{python}
for etapes in [3, 6, 10]:
    print(etapes, round(0.90 ** etapes * 100), '%')
\end{python}
\begin{console}
3 73 %
6 53 %
10 35 %
\end{console}
Même avec d'excellents rendements par étape, une synthèse longue gaspille beaucoup de matière. La chimie
\res{\text{plus durable}} cherche des synthèses courtes, à forte économie d'atomes, utilisant des catalyseurs, des solvants peu
dangereux et des matières premières renouvelables.
\editionRetenir{Le rendement global d'une synthèse en plusieurs étapes est le produit des rendements de chaque étape.
L'économie d'atomes mesure, dès l'écriture des équations, la part des réactifs qui finit dans le produit.}
\end{blocEdition}

\editionPage{defi-cor-10-1}{D10 · Corrigé}{Le meilleur carburant ?}{\hyperref[defi-ex-10]{Énoncé p.~\pageref*{defi-ex-10}} · Démarche et vérifications}
\begin{blocEdition}{defi-cor-10-1}
\editionEtape{1. Les pouvoirs calorifiques}
Dihydrogène : $\dfrac{242}{2{,}0}=\res{\qty{121}{MJ.kg^{-1}}}$ ; méthane : $\dfrac{802}{16{,}0}\approx\res{\qty{50,1}{MJ.kg^{-1}}}$ ; éthanol : $\dfrac{1235}{46{,}0}\approx\res{\qty{26,8}{MJ.kg^{-1}}}$.
\editionEtape{2. L'énergie par litre}
Dihydrogène : $121\times0{,}071\approx\res{\qty{8,6}{MJ.L^{-1}}}$ ; méthane : $50{,}1\times0{,}42\approx\res{\qty{21,0}{MJ.L^{-1}}}$ ; kérosène : $43{,}0\times0{,}80=\res{\qty{34,4}{MJ.L^{-1}}}$ ;
éthanol : $26{,}8\times0{,}789\approx\res{\qty{21,1}{MJ.L^{-1}}}$. Par kilogramme : dihydrogène, méthane, kérosène, éthanol. Par litre : kérosène,
éthanol, méthane, dihydrogène. Les deux classements sont presque \res{\text{inversés}}.
\editionEtape{3. Pour un gigajoule}
Dihydrogène : $\dfrac{1000}{121}\approx\res{\qty{8,3}{kg}}$, occupant $\dfrac{1000}{8{,}6}\approx\res{\qty{116}{L}}$. Kérosène : $\dfrac{1000}{43{,}0}\approx\res{\qty{23,3}{kg}}$, occupant
$\dfrac{1000}{34{,}4}\approx\res{\qty{29,1}{L}}$. Le dihydrogène est près de trois fois plus léger, mais quatre fois plus encombrant.
\editionEtape{4. Le dioxyde de carbone}
$\ce{2C12H26 + 37O2 -> 24CO2 + 26H2O}$. Avec $M(\ce{C12H26})=\qty{170,0}{g.mol^{-1}}$, $|E_\mathrm{r}|=43{,}0\times170{,}0=\qty{7310}{kJ.mol^{-1}}$ ; une mole libère
$12\times44{,}0=\qty{528}{g}$ de $\ce{CO2}$, soit $\dfrac{528}{7{,}31}\approx\res{\qty{72,2}{g.MJ^{-1}}}$, contre \qty{54,9}{g.MJ^{-1}} pour le méthane, \qty{71,3}{g.MJ^{-1}} pour l'éthanol et
zéro pour le dihydrogène.
\end{blocEdition}

\editionPage{defi-cor-10-2}{D10 · Corrigé (suite)}{Le meilleur carburant ?}{\hyperref[defi-ex-10]{Énoncé p.~\pageref*{defi-ex-10}} · Démarche et vérifications}
\begin{blocEdition}{defi-cor-10-2}
\label{figedition-defi-cor-10-2-1}\schema[corrige]{defi-10-cor}
\editionEtape{5. Les étages d'une fusée}
Au décollage, il faut une poussée énorme pendant peu de temps, avec des réservoirs compacts : le kérosène, très
énergétique \res{\text{par litre}}, limite la taille et la masse des réservoirs. Pour les étages supérieurs, chaque kilogramme
emporté coûte très cher : le dihydrogène, imbattable \res{\text{par kilogramme}}, s'impose.
\editionEtape{6. Un programme de classement}
\begin{python}
carburants = {'H2': [121, 0.071], 'CH4': [50.1, 0.42],
              'kerosene': [43.0, 0.80], 'ethanol': [26.8, 0.789]}
par_litre = {nom: round(pc * rho, 1) for nom, (pc, rho) in carburants.items()}
print(sorted(par_litre, key=par_litre.get, reverse=True))
print(par_litre)
\end{python}
\begin{console}
['kerosene', 'ethanol', 'CH4', 'H2']
{'H2': 8.6, 'CH4': 21.0, 'kerosene': 34.4, 'ethanol': 21.1}
\end{console}
\editionRetenir{Un carburant se juge sur plusieurs critères : énergie par kilogramme, énergie par litre, émissions de
dioxyde de carbone, mais aussi mode de production et sécurité. Aucun n'est le meilleur sur tous les plans.}
\end{blocEdition}

\editionPage{defi-cor-11-1}{D11 · Corrigé}{Ce qui tient le noyau}{\hyperref[defi-ex-11]{Énoncé p.~\pageref*{defi-ex-11}} · Démarche et vérifications}
\begin{blocEdition}{defi-cor-11-1}
\editionEtape{1. La répulsion électrostatique}
$F_\mathrm{e}=k\,\dfrac{e^2}{d^2}=\num{9,0e9}\times\dfrac{(\num{1,60e-19})^2}{(\num{2,0e-15})^2}=\res{\qty{57,6}{N}}$. Les deux charges sont positives : la force est \res{\text{répulsive}}.
\editionEtape{2. Une force à notre échelle}
$m=\dfrac{F_\mathrm{e}}{g}=\dfrac{57{,}6}{9{,}81}\approx\res{\qty{5,9}{kg}}$ : chaque proton est repoussé aussi fort qu'un sac de près de six kilogrammes est
attiré par la Terre, alors que sa masse ne vaut que \qty{1,67e-27}{kg} !
\editionEtape{3. L'attraction gravitationnelle}
$F_\mathrm{g}=G\,\dfrac{m_\mathrm{p}^2}{d^2}=\num{6,67e-11}\times\dfrac{(\num{1,67e-27})^2}{(\num{2,0e-15})^2}=\res{\qty{4,65e-35}{N}}$. Le poids d'un proton vaut
$m_\mathrm{p}\,g=\qty{1,64e-26}{N}$ : l'attraction entre les deux protons est encore $\num{3,5e8}$ fois plus faible que le poids d'un proton.
\editionEtape{4. Le rapport des deux forces}
$\dfrac{F_\mathrm{e}}{F_\mathrm{g}}=\dfrac{57{,}6}{\num{4,65e-35}}\approx\res{\num{1,24e36}}$. La gravitation est \res{\text{totalement incapable}} de retenir les protons.
\editionEtape{5. La distance pour un newton}
$d=e\,\sqrt{\dfrac{k}{F}}=\num{1,60e-19}\times\sqrt{\num{9,0e9}}\approx\res{\qty{1,52e-14}{m}}$, soit environ 7,6 fois la distance initiale : la répulsion
diminue vite, mais reste considérable à l'échelle d'une particule.
\label{figedition-defi-cor-11-1-1}\schema[corrige]{defi-11-cor}
\end{blocEdition}

\editionPage{defi-cor-11-2}{D11 · Corrigé (suite)}{Ce qui tient le noyau}{\hyperref[defi-ex-11]{Énoncé p.~\pageref*{defi-ex-11}} · Démarche et vérifications}
\begin{blocEdition}{defi-cor-11-2}
\editionEtape{6. Un rapport indépendant de la distance}
\begin{python}
k, G, e, mp = 9.0e9, 6.67e-11, 1.60e-19, 1.67e-27
for d in [2.0e-15, 1.0e-10, 1.0]:
    Fe = k * e**2 / d**2
    Fg = G * mp**2 / d**2
    print('d =', d, 'm : Fe/Fg =', f'{Fe / Fg:.3g}')
\end{python}
\begin{console}
d = 2e-15 m : Fe/Fg = 1.24e+36
d = 1e-10 m : Fe/Fg = 1.24e+36
d = 1.0 m : Fe/Fg = 1.24e+36
\end{console}
Le rapport vaut $\dfrac{F_\mathrm{e}}{F_\mathrm{g}}=\dfrac{k\,e^2}{G\,m_\mathrm{p}^2}$ : le facteur $d^2$ se \res{\text{simplifie}}, car les deux forces varient en $\dfrac{1}{d^2}$. Quelle que soit la
distance, la répulsion électrostatique entre deux protons l'emporte de $10^{36}$ sur leur attraction gravitationnelle.
\editionEtape{7. Une troisième interaction}
Puisque le noyau existe, une autre interaction attire les nucléons entre eux : l'\res{\text{interaction forte}}. Elle doit être
\res{\text{attractive}}, bien plus intense que la répulsion électrostatique à la distance de \qty{e-15}{m}, et de \res{\text{très courte portée}} :
on ne la ressent pas à l'échelle de l'atome, où la répulsion entre noyaux domine.
\editionRetenir{Quatre interactions fondamentales régissent l'Univers : gravitationnelle, électromagnétique, forte et faible.
Chacune domine à sa propre échelle : la forte dans le noyau, l'électromagnétique dans l'atome et la matière, la
gravitation entre les astres.}
\end{blocEdition}

\editionPage{defi-cor-12-1}{D12 · Corrigé}{Les hémisphères de Magdebourg}{\hyperref[defi-ex-12]{Énoncé p.~\pageref*{defi-ex-12}} · Démarche et vérifications}
\begin{blocEdition}{defi-cor-12-1}
\editionEtape{1. Avec un vide parfait}
$S=\pi\,R^2=\pi\times0{,}25^2\approx\qty{0,196}{m^2}$ et $F=P_\mathrm{atm}\,S=\num{1,013e5}\times0{,}196\approx\res{\qty{1,99e4}{N}}$ : le poids d'une masse de deux tonnes.
\editionEtape{2. Les chevaux et l'arbre}
Chaque attelage doit tirer avec \qty{1,99e4}{N}, soit $\dfrac{\num{1,99e4}}{8}\approx\res{\qty{2,5e3}{N}}$ par cheval, plus que ce qu'un cheval peut fournir
durablement. Si un seul attelage tire, le crochet fixé à l'arbre exerce une force opposée de même valeur : le second
attelage ne fait que \res{\text{remplacer l'arbre}}. Pour doubler l'effort, il aurait fallu atteler les seize chevaux du même côté.
\editionEtape{3. Un coup de pompe}
L'air intérieur, à la pression $P$ sous le volume $V$, se répartit dans $V+v$ à température constante. D'après la loi de Mariotte,
$P\,V=P'\,(V+v)$, donc $P'=P\times\res{\dfrac{V}{V+v}}=P\times\dfrac{65{,}4}{67{,}4}\approx\num{0,970}\,P$ : chaque coup retire \qty{3}{\percent} de l'air restant.
\editionEtape{4. Compter les coups}
\begin{python}
V, v = 65.4, 2.0          # L : spheres, cylindre de la pompe
P, n = 1.013, 0           # bar, nombre de coups
while P > 0.10:
    P = P * V / (V + v)   # loi de Mariotte
    n = n + 1
print(n, 'coups :', round(P, 4), 'bar')
\end{python}
\begin{console}
77 coups : 0.0996 bar
\end{console}
Il faut \res{\text{77 coups}} de pompe : la pression décroît de moins en moins vite, car chaque coup retire une fraction de l'air restant.
\end{blocEdition}

\editionPage{defi-cor-12-2}{D12 · Corrigé (suite)}{Les hémisphères de Magdebourg}{\hyperref[defi-ex-12]{Énoncé p.~\pageref*{defi-ex-12}} · Démarche et vérifications}
\begin{blocEdition}{defi-cor-12-2}
\label{figedition-defi-cor-12-2-1}\schema[corrige]{defi-12-cor}
\editionEtape{5. Avec la pression résiduelle}
$F=(P_\mathrm{atm}-P)\,S=(\num{1,013e5}-\num{0,10e5})\times0{,}196\approx\res{\qty{1,79e4}{N}}$ : à peine \qty{10}{\percent} de moins qu'avec un vide parfait. Un vide
imparfait suffit donc à rendre la séparation impossible.
\editionEtape{6. La ventouse}
$F_\mathrm{max}=P_\mathrm{atm}\,\pi\,r^2=\num{1,013e5}\times\pi\times0{,}050^2\approx\res{\qty{796}{N}}$, le poids d'une masse de \qty{80}{kg}. En pratique, il reste de
l'air sous la ventouse et de l'air s'infiltre peu à peu : la différence de pression, et donc la force, sont bien plus faibles.
\editionRetenir{Ce n'est pas le vide qui « aspire » : c'est l'air extérieur qui pousse. Une différence de pression $\Delta P$ sur une surface $S$
produit une force $\Delta P\times S$, considérable dès que la surface atteint quelques décimètres carrés.}
\end{blocEdition}

\editionPage{defi-cor-13-1}{D13 · Corrigé}{Freiner à temps}{\hyperref[defi-ex-13]{Énoncé p.~\pageref*{defi-ex-13}} · Démarche et vérifications}
\begin{blocEdition}{defi-cor-13-1}
\editionEtape{1. Le temps de réaction}
Le conducteur n'a pas encore freiné : mouvement \res{\text{rectiligne uniforme}}. $d=v\times\qty{1,0}{s}$ : \res{\qty{13,9}{m}} à \qty{50}{km/h}
($\dfrac{50}{3{,}6}\approx\qty{13,9}{m.s^{-1}}$), \res{\qty{25,0}{m}} à \qty{90}{km/h} et \res{\qty{36,1}{m}} à \qty{130}{km/h}.
\editionEtape{2. Le freinage}
$\sum F=0{,}8\,m\,g=0{,}8\times1200\times9{,}81\approx\res{\qty{9,42e3}{N}}$, opposée à la vitesse. Toutes les \qty{0,01}{s},
$\Delta v=\dfrac{\sum F}{m}\,\Delta t=0{,}8\times9{,}81\times0{,}01\approx\res{\qty{0,078}{m.s^{-1}}}$. La masse se \res{\text{simplifie}} : une voiture lourde ne freine
pas moins bien, à adhérence égale.
\editionEtape{3. La simulation}
\begin{python}
v, x, dt = 130 / 3.6, 0, 0.01      # m/s, m, s
while v > 0:
    x = x + v * dt                 # la voiture avance
    v = v - 0.8 * 9.81 * dt        # Delta v = (F / m) x dt
print(round(x, 1), 'm')
\end{python}
\begin{console}
83.3 m
\end{console}
\editionEtape{4. Les distances d'arrêt}
Distance d'arrêt = distance de réaction + distance de freinage : \res{\qty{26}{m}} et \res{\qty{34}{m}} à \qty{50}{km/h} ; \res{\qty{65}{m}} et \res{\qty{89}{m}}
à \qty{90}{km/h} ; \res{\qty{119}{m}} et \res{\qty{169}{m}} à \qty{130}{km/h} (route sèche puis mouillée).
\end{blocEdition}

\editionPage{defi-cor-13-2}{D13 · Corrigé (suite)}{Freiner à temps}{\hyperref[defi-ex-13]{Énoncé p.~\pageref*{defi-ex-13}} · Démarche et vérifications}
\begin{blocEdition}{defi-cor-13-2}
\label{figedition-defi-cor-13-2-1}\schema[corrige]{defi-13-cor}
\editionEtape{5. Quand la vitesse double}
$\dfrac{49{,}3}{12{,}4}\approx\res{4{,}0}$ : la distance de freinage est \res{\text{multipliée par quatre}}. En effet, il faut deux fois plus de temps pour
annuler une vitesse double, et la voiture roule en moyenne deux fois plus vite pendant ce temps.
\editionEtape{6. La règle des deux secondes}
À \qty{130}{km/h}, deux secondes représentent $36{,}1\times2\approx\res{\qty{72}{m}}$, moins que la distance d'arrêt de \qty{119}{m}. Mais le véhicule de
devant freine lui aussi : il faut surtout couvrir le \res{\text{temps de réaction}} (\qty{36}{m}) avec une marge. La règle devient dangereuse
sur route mouillée, ou si le véhicule de devant s'arrête brutalement contre un obstacle.
\editionRetenir{La distance de freinage est proportionnelle au carré de la vitesse : rouler à \qty{130}{km/h} au lieu de \qty{90}{km/h} fait plus que
doubler la distance de freinage. D'où les limitations de vitesse plus basses par temps de pluie.}
\end{blocEdition}

\editionPage{defi-cor-14-1}{D14 · Corrigé}{Ne pas faire sauter les plombs}{\hyperref[defi-ex-14]{Énoncé p.~\pageref*{defi-ex-14}} · Démarche et vérifications}
\begin{blocEdition}{defi-cor-14-1}
\editionEtape{1. Les intensités}
$I=\dfrac{P}{U}$ : plaque \res{\qty{15,7}{A}}, four \res{\qty{10,9}{A}}, bouilloire \res{\qty{9,57}{A}}, lave-linge \res{\qty{8,70}{A}}, réfrigérateur \qty{0,65}{A}, éclairage
\qty{0,43}{A}. La plaque, l'appareil le plus puissant, est protégée à \res{\qty{32}{A}} : avec un calibre de \qty{16}{A}, elle serait à la limite de
la coupure (et son câble, plus gros, supporte davantage).
\editionEtape{2. Un circuit de prises}
$9{,}57+8{,}70\approx\res{\qty{18,3}{A}}>\qty{16}{A}$ : le disjoncteur du circuit de prises \res{\text{coupe}} ; il faut les utiliser l'un après l'autre.
\editionEtape{3. L'abonnement}
$I_\mathrm{max}=\dfrac{6000}{230}\approx\res{\qty{26,1}{A}}$. Four, plaque et bouilloire appellent $2500+3600+2200=\qty{8300}{W}>\qty{6000}{W}$ : le disjoncteur général
\res{\text{coupe toute l'installation}}.
\editionEtape{4. Toutes les combinaisons}
\begin{python}
from itertools import combinations
appareils = {'bouilloire': 2200, 'four': 2500, 'plaque': 3600,
             'lave-linge': 2000, 'frigo': 150, 'lumiere': 100}
n = 0
for k in range(1, 7):
    for groupe in combinations(appareils, k):
        if sum(appareils[a] for a in groupe) > 6000:
            n = n + 1
print(n, 'combinaisons sur 63 font disjoncter')
\end{python}
\begin{console}
25 combinaisons sur 63 font disjoncter
\end{console}
\end{blocEdition}

\editionPage{defi-cor-14-2}{D14 · Corrigé (suite)}{Ne pas faire sauter les plombs}{\hyperref[defi-ex-14]{Énoncé p.~\pageref*{defi-ex-14}} · Démarche et vérifications}
\begin{blocEdition}{defi-cor-14-2}
\label{figedition-defi-cor-14-2-1}\schema[corrige]{defi-14-cor}
\editionEtape{5. La rallonge}
$R=\num{1,7e-8}\times\dfrac{50}{\num{1,5e-6}}\approx\res{\qty{0,567}{\ohm}}$. Avec la bouilloire, $P=R\,I^2=0{,}567\times9{,}57^2\approx\res{\qty{52}{W}}$ dissipés dans la rallonge, et
une chute de tension $R\,I\approx\qty{5,4}{V}$. Avec un câble de \qty{2,5}{mm^2}, les pertes tomberaient à \qty{31}{W}.
\editionEtape{6. Dérouler l'enrouleur}
Enroulé, le câble ne peut pas évacuer ces \qty{52}{W} : comme un petit radiateur enfermé, il s'échauffe jusqu'à faire fondre son isolant.
Déroulé, il se refroidit dans l'air sur toute sa longueur.
\editionRetenir{Un disjoncteur protège les câbles contre l'échauffement par effet Joule : l'intensité dans un circuit ne doit pas dépasser
son calibre. Les intensités des appareils branchés sur un même circuit s'additionnent.}
\end{blocEdition}

\editionPage{defi-cor-15-1}{D15 · Corrigé}{Rouler moins vite pour consommer moins}{\hyperref[defi-ex-15]{Énoncé p.~\pageref*{defi-ex-15}} · Démarche et vérifications}
\begin{blocEdition}{defi-cor-15-1}
\editionEtape{1. Le bilan à vitesse constante}
$\Delta E_\mathrm{c}=\res{0}$ : d'après le théorème de l'énergie cinétique, la somme des travaux est nulle. Le poids et la réaction de la route ne
travaillent pas ; le travail \res{\text{moteur}} compense exactement le travail \res{\text{résistant}} des frottements.
\editionEtape{2. Les frottements de l'air}
$f=\dfrac{1}{2}\times1{,}2\times0{,}65\times v^2$ : \res{\qty{244}{N}} à \qty{25,0}{m.s^{-1}} (\qty{90}{km/h}), \res{\qty{364}{N}} à \qty{30,6}{m.s^{-1}} et \res{\qty{509}{N}} à \qty{36,1}{m.s^{-1}}. Ils
font plus que doubler entre 90 et \qty{130}{km/h}.
\editionEtape{3. À 130 km/h}
$W_\mathrm{moteur}=(509+150)\times\num{1,00e5}\approx\res{\qty{6,59e7}{J}}$. Chaque litre fournit $0{,}30\times\num{35e6}=\qty{10,5e6}{J}$ : $\dfrac{\num{6,59e7}}{\num{10,5e6}}\approx\res{\qty{6,27}{L}}$
pour 100 km.
\editionEtape{4. Le programme}
\begin{python}
Cs, rho, fr = 0.65, 1.2, 150        # m2, kg/m3, N
for v_kmh in [90, 110, 130]:
    v = v_kmh / 3.6
    f = 0.5 * rho * Cs * v**2 + fr  # forces resistantes (N)
    W = f * 100e3                   # travail moteur sur 100 km (J)
    print(v_kmh, 'km/h :', round(W / 10.5e6, 2), 'L/100 km')
\end{python}
\begin{console}
90 km/h : 3.75 L/100 km
110 km/h : 4.9 L/100 km
130 km/h : 6.27 L/100 km
\end{console}
\end{blocEdition}

\editionPage{defi-cor-15-2}{D15 · Corrigé (suite)}{Rouler moins vite pour consommer moins}{\hyperref[defi-ex-15]{Énoncé p.~\pageref*{defi-ex-15}} · Démarche et vérifications}
\begin{blocEdition}{defi-cor-15-2}
\label{figedition-defi-cor-15-2-1}\schema[corrige]{defi-15-cor}
\editionEtape{5. L'économie}
$6{,}27-4{,}90\approx\res{\qty{1,37}{L}}$ pour 100 km, soit \res{\qty{22}{\percent}} de moins. Dioxyde de carbone : $6{,}27\times2{,}3\approx\qty{14,4}{kg}$ contre
$4{,}90\times2{,}3\approx\qty{11,3}{kg}$, environ \res{\qty{3}{kg}} de moins tous les 100 km.
\editionEtape{6. Le temps perdu}
100 km prennent \qty{46}{min} à \qty{130}{km/h} et \qty{55}{min} à \qty{110}{km/h} : on perd un peu plus de \res{\qty{8}{min}} par 100 km pour économiser plus d'un litre.
Le modèle néglige les accélérations et la consommation des équipements, mais l'ordre de grandeur est réaliste.
\editionRetenir{À vitesse constante, le moteur ne sert qu'à compenser le travail des frottements. Comme les frottements de l'air
croissent avec le carré de la vitesse, rouler un peu moins vite fait beaucoup baisser la consommation.}
\end{blocEdition}

\editionPage{defi-cor-16-1}{D16 · Corrigé}{Stocker l'électricité dans un lac}{\hyperref[defi-ex-16]{Énoncé p.~\pageref*{defi-ex-16}} · Démarche et vérifications}
\begin{blocEdition}{defi-cor-16-1}
\editionEtape{1. L'énergie potentielle stockée}
$m=\rho\,V=1000\times\num{1,0e8}=\qty{1,0e11}{kg}$, d'où $E_\mathrm{p}=m\,g\,h=\num{1,0e11}\times9{,}81\times900\approx\res{\qty{8,83e14}{J}}$, soit $\dfrac{\num{8,83e14}}{\num{3,6e12}}\approx\res{\qty{245}{GWh}}$.
\editionEtape{2. Le rendement du stockage}
Énergie restituée : $0{,}90\times245\approx\res{\qty{221}{GWh}}$. Énergie de pompage : $\dfrac{245}{0{,}85}\approx\res{\qty{289}{GWh}}$. Rendement global :
$\dfrac{221}{289}=0{,}85\times0{,}90\approx\res{\num{0,765}}$ : on récupère les trois quarts de l'électricité stockée.
\editionEtape{3. Le bilan énergétique}
\label{figedition-defi-cor-16-1-1}\schema[corrige]{defi-16-cor}
\end{blocEdition}

\editionPage{defi-cor-16-2}{D16 · Corrigé (suite)}{Stocker l'électricité dans un lac}{\hyperref[defi-ex-16]{Énoncé p.~\pageref*{defi-ex-16}} · Démarche et vérifications}
\begin{blocEdition}{defi-cor-16-2}
\editionEtape{4. Le programme}
\begin{python}
rho, g, h = 1000, 9.81, 900          # kg/m3, N/kg, m
for V in [1e6, 1e7, 1e8]:            # volumes en m3
    E = rho * V * g * h * 0.90 / 3.6e9   # energie restituee en MWh
    print(round(V), 'm3 :', round(E), 'MWh')
\end{python}
\begin{console}
1000000 m3 : 2207 MWh
10000000 m3 : 22072 MWh
100000000 m3 : 220725 MWh
\end{console}
L'énergie stockée est proportionnelle au volume d'eau (et au dénivelé).
\editionEtape{5. Le débit et l'autonomie}
Un débit de \qty{1}{m^3.s^{-1}} fournit $0{,}90\times1000\times9{,}81\times900\approx\qty{7,95}{MW}$. Pour \qty{1800}{MW} : $Q=\dfrac{1800}{7{,}95}\approx\res{\qty{226}{m^3.s^{-1}}}$. La réserve
dure $\dfrac{\num{1,0e8}}{226}\approx\qty{4,4e5}{s}$, soit environ \res{\qty{123}{h}}, plus de cinq jours.
\editionEtape{6. Comparaison avec des batteries}
$\dfrac{\num{221e6}}{52}\approx\res{\num{4,2e6}}$ batteries de voiture : plus de quatre millions de véhicules. Les STEP restent, de loin, le moyen le plus
efficace de stocker de grandes quantités d'électricité, mais il faut deux lacs et un fort dénivelé.
\editionRetenir{L'énergie potentielle de pesanteur permet de stocker de l'énergie : on la convertit en électricité au moment voulu.
Chaque conversion a un rendement inférieur à 1, et les pertes se multiplient.}
\end{blocEdition}

\editionPage{defi-cor-17-1}{D17 · Corrigé}{Localiser l'épicentre d'un séisme}{\hyperref[defi-ex-17]{Énoncé p.~\pageref*{defi-ex-17}} · Démarche et vérifications}
\begin{blocEdition}{defi-cor-17-1}
\editionEtape{1. Distance et retard}
Les ondes P arrivent après $\dfrac{d}{6{,}0}$, les ondes S après $\dfrac{d}{3{,}5}$ : $\Delta t=d\left(\dfrac{1}{3{,}5}-\dfrac{1}{6{,}0}\right)=k\,d$, avec $k\approx\qty{0,119}{s/km}$. Donc $\res{d=\dfrac{\Delta t}{k}}$.
\editionEtape{2. Les trois distances}
$d_\mathrm{A}=\dfrac{7{,}46}{\num{0,119}}\approx\res{\qty{62,7}{km}}$ ; $d_\mathrm{B}\approx\res{\qty{54,1}{km}}$ ; $d_\mathrm{C}\approx\res{\qty{74,3}{km}}$.
\editionEtape{3. Les cercles}
\label{figedition-defi-cor-17-1-1}\schema[corrige]{defi-17}
\end{blocEdition}

\editionPage{defi-cor-17-2}{D17 · Corrigé (suite)}{Localiser l'épicentre d'un séisme}{\hyperref[defi-ex-17]{Énoncé p.~\pageref*{defi-ex-17}} · Démarche et vérifications}
\begin{blocEdition}{defi-cor-17-2}
\editionEtape{4. Une recherche systématique}
\begin{python}
stations = [(0, 0, 7.46), (100, 0, 6.44), (0, 80, 8.85)]   # km, km, s
k = 1 / 3.5 - 1 / 6.0
meilleur = None
for x in range(0, 101):
    for y in range(0, 101):
        e = 0
        for (xs, ys, dt) in stations:
            d = ((x - xs)**2 + (y - ys)**2) ** 0.5
            e = e + (d - dt / k) ** 2
        if meilleur is None or e < meilleur[0]:
            meilleur = (e, x, y)
print('epicentre :', meilleur[1], 'km ;', meilleur[2], 'km')
\end{python}
\begin{console}
epicentre : 55 km ; 30 km
\end{console}
\editionEtape{5. L'arrivée des ondes P}
En A, $t=\dfrac{62{,}7}{6{,}0}\approx\res{\qty{10,4}{s}}$ après le début du séisme.
\editionEtape{6. Pourquoi trois stations ?}
Deux cercles se coupent en général en \res{\text{deux points}} : la troisième station permet de choisir le bon. En pratique, on utilise
de nombreuses stations, ce qui améliore la précision.
\editionRetenir{Mesurer un retard entre deux ondes de célérités différentes donne une distance, sans connaître la date de départ.
Trois distances suffisent pour localiser une source dans un plan.}
\end{blocEdition}

\editionPage{defi-cor-18-1}{D18 · Corrigé}{Photographier la Lune}{\hyperref[defi-ex-18]{Énoncé p.~\pageref*{defi-ex-18}} · Démarche et vérifications}
\begin{blocEdition}{defi-cor-18-1}
\editionEtape{1. Une image au foyer}
$\dfrac{1}{\alg{OA'}}=\dfrac{1}{f'}+\dfrac{1}{\alg{OA}}$ avec $|\alg{OA}|=\qty{3,844e8}{m}$ : le second terme est des milliards de fois plus petit que le premier, donc
$\res{\alg{OA'}\approx f'}$ : l'image se forme dans le plan focal, où l'on place le capteur.
\editionEtape{2. Le diamètre apparent}
Le rayon issu d'un bord de la Lune et passant par O n'est pas dévié. Par Thalès, $\dfrac{d'}{D}=\dfrac{\alg{OA'}}{|\alg{OA}|}\approx\dfrac{f'}{d}$, donc $\res{d'=f'\times\dfrac{D}{d}=f'\times\alpha}$
avec $\alpha=\dfrac{3474}{384400}\approx\res{\qty{9,04e-3}{rad}}$, soit environ \ang{0,52} : la Lune paraît aussi grande qu'une pièce de 1 cm vue à 1 m.
\editionEtape{3. Avec un téléobjectif}
$d'=300\times\num{9,04e-3}\approx\res{\qty{2,71}{mm}}$, environ \res{\qty{11}{\percent}} des \qty{24}{mm} de hauteur.
\editionEtape{4. Avec le smartphone}
$d'=4\times\num{9,04e-3}\approx\qty{0,036}{mm}=\res{\qty{36}{\micro m}}$, soit $\dfrac{36}{1{,}4}\approx\res{26}$ pixels : un petit disque sans détail.
\label{figedition-defi-cor-18-1-1}\schema[corrige]{defi-18-cor}
\end{blocEdition}

\editionPage{defi-cor-18-2}{D18 · Corrigé (suite)}{Photographier la Lune}{\hyperref[defi-ex-18]{Énoncé p.~\pageref*{defi-ex-18}} · Démarche et vérifications}
\begin{blocEdition}{defi-cor-18-2}
\editionEtape{5. Le programme}
\begin{python}
alpha = 3474 / 384400              # diametre apparent (rad)
for fp in [4, 50, 300, 1200]:      # focales en mm
    print(fp, 'mm :', round(fp * alpha, 2), 'mm')
\end{python}
\begin{console}
4 mm : 0.04 mm
50 mm : 0.45 mm
300 mm : 2.71 mm
1200 mm : 10.84 mm
\end{console}
\editionEtape{6. Remplir le capteur}
$f'=\dfrac{24}{\num{9,04e-3}}\approx\res{\qty{2656}{mm}}$, plus de \qty{2,6}{m}. Seuls des instruments de très grande focale (télescopes) donnent une image assez
grande de la Lune ou des planètes ; leur grand diamètre recueille aussi davantage de lumière.
\editionRetenir{Pour un objet très éloigné, la taille de l'image est proportionnelle à la distance focale : $d'=f'\times\alpha$. Un objectif de
smartphone, de très courte focale, donne une toute petite image des astres.}
\end{blocEdition}

\editionPage{defi-cor-19-1}{D19 · Corrigé}{Pourquoi le ciel est-il bleu ?}{\hyperref[defi-ex-19]{Énoncé p.~\pageref*{defi-ex-19}} · Démarche et vérifications}
\begin{blocEdition}{defi-cor-19-1}
\editionEtape{1. Bleu contre rouge}
$\left(\dfrac{700}{450}\right)^4\approx\res{5{,}86}$ : le bleu est près de six fois plus diffusé que le rouge ; $\left(\dfrac{550}{450}\right)^4\approx\res{2{,}23}$ fois plus que le vert.
\editionEtape{2. Le programme}
\begin{python}
for lam in [400, 450, 500, 550, 600, 650, 700]:   # nm
    print(lam, 'nm :', round((700 / lam) ** 4, 2))
\end{python}
\begin{console}
400 nm : 9.38
450 nm : 5.86
500 nm : 3.84
550 nm : 2.62
600 nm : 1.85
650 nm : 1.35
700 nm : 1.0
\end{console}
\label{figedition-defi-cor-19-1-1}\schema[corrige]{defi-19-cor}
\end{blocEdition}

\editionPage{defi-cor-19-2}{D19 · Corrigé (suite)}{Pourquoi le ciel est-il bleu ?}{\hyperref[defi-ex-19]{Énoncé p.~\pageref*{defi-ex-19}} · Démarche et vérifications}
\begin{blocEdition}{defi-cor-19-2}
\editionEtape{3. Un ciel bleu}
Quand on regarde le ciel loin du Soleil, on reçoit uniquement la lumière \res{\text{diffusée}} par l'air, qui vient de toutes les directions. Elle
est enrichie en courtes longueurs d'onde : elle paraît \res{\text{bleue}}.
\editionEtape{4. Pourquoi pas violet ?}
Le Soleil émet moins de violet que de bleu, une partie du violet est absorbée en haute altitude, et l'œil y est peu sensible : les cônes
S ont leur maximum vers \qty{420}{nm} et les cônes M et L sont aussi un peu excités. Le mélange perçu est un \res{\text{bleu clair}}.
\editionEtape{5. Le Soleil couchant}
$\dfrac{1}{\cos\ang{85}}\approx\res{11{,}5}$ : la lumière traverse onze fois plus d'air. Tout au long de ce trajet, le bleu est diffusé hors du faisceau ; la lumière
directe qui nous parvient s'est appauvrie en bleu : le Soleil paraît \res{\text{orangé puis rouge}}.
\editionEtape{6. Des nuages blancs}
Les gouttelettes diffusent \res{\text{toutes les radiations}} de la même façon : la lumière diffusée garde la composition de la lumière solaire et paraît
blanche (ou grise, si le nuage est épais et arrête la lumière).
\editionRetenir{La couleur du ciel est celle de la lumière diffusée ; la couleur du Soleil, celle de la lumière transmise. Les deux sont
complémentaires : ce que le ciel diffuse en bleu manque à la lumière directe.}
\end{blocEdition}

\editionPage{defi-cor-20-1}{D20 · Corrigé}{De quoi les étoiles sont-elles faites ?}{\hyperref[defi-ex-20]{Énoncé p.~\pageref*{defi-ex-20}} · Démarche et vérifications}
\begin{blocEdition}{defi-cor-20-1}
\editionEtape{1. Un spectre d'absorption}
Le fond est un spectre continu barré de raies noires : c'est un spectre d'\res{\text{absorption}}. La surface très chaude de l'étoile émet
toutes les radiations ; les atomes de son atmosphère absorbent celles qui correspondent exactement à l'écart entre deux de leurs niveaux.
\editionEtape{2. Une raie de l'hydrogène}
$E_5-E_2=-\dfrac{13{,}6}{25}+\dfrac{13{,}6}{4}=-0{,}544+3{,}40=\qty{2,86}{eV}$, d'où $\lambda=\dfrac{\num{6,63e-34}\times\num{3,00e8}}{2{,}86\times\num{1,60e-19}}=\res{\qty{435}{nm}}$. L'écart
relatif avec \qty{434,0}{nm} vaut \qty{0,3}{\percent} : il vient de l'\res{\text{arrondi des constantes}} $h$, $c$ et \qty{13,6}{eV}.
\editionEtape{3. Le programme}
\begin{python}
mesures = [656.3, 589.0, 527.0, 517.3, 486.1, 434.0, 393.4]   # nm
tables = {'H': [656.3, 486.1, 434.0, 410.2], 'Na': [589.0, 589.6],
          'Mg': [517.3, 518.4], 'Fe': [527.0, 438.4], 'Ca': [393.4, 396.8]}
trouves = []
for lam in mesures:
    for el in tables:
        for r in tables[el]:
            if abs(lam - r) < 0.5 and el not in trouves:
                trouves.append(el)
print(trouves)
\end{python}
\begin{console}
['H', 'Na', 'Fe', 'Mg', 'Ca']
\end{console}
Le test \code{el not in trouves} évite les doublons : l'hydrogène, qui signe trois raies, n'est ajouté qu'une fois.
\label{figedition-defi-cor-20-1-1}\schema[corrige]{defi-20}
\end{blocEdition}

\editionPage{defi-cor-20-2}{D20 · Corrigé (suite)}{De quoi les étoiles sont-elles faites ?}{\hyperref[defi-ex-20]{Énoncé p.~\pageref*{defi-ex-20}} · Démarche et vérifications}
\begin{blocEdition}{defi-cor-20-2}
\editionEtape{4. La raie du calcium}
$E=\dfrac{h\,c}{\lambda}=\dfrac{\num{6,63e-34}\times\num{3,00e8}}{\num{393,4e-9}}=\qty{5,06e-19}{J}$, soit \res{\qty{3,16}{eV}}. Avec $\lambda<\qty{400}{nm}$, la raie est dans le
\res{\text{proche ultraviolet}}, à la limite du violet.
\editionEtape{5. Une atmosphère chaude}
Les raies visibles de l'hydrogène correspondent à des transitions depuis le niveau $n=2$. À température ambiante, les atomes sont
tous dans l'état fondamental et n'absorbent que dans l'UV. Si l'étoile absorbe ces raies, c'est qu'une partie de ses atomes
d'hydrogène est déjà sur le niveau 2 : son atmosphère est \res{\text{très chaude}} (plusieurs milliers de degrés).
\editionEtape{6. Un élément inconnu}
$589{,}0-587{,}6=\qty{1,4}{nm}>\qty{0,5}{nm}$ : le programme \res{\text{ne l'attribue pas au sodium}}. Aucun élément connu n'émettant à cette longueur
d'onde, on en a déduit l'existence d'un \res{\text{nouvel élément}}, baptisé \textbf{hélium} (du grec \emph{hélios}, le Soleil). Il ne fut isolé sur
Terre qu'en 1895.
\editionRetenir{Chaque élément possède ses propres niveaux d'énergie, donc ses propres raies : un spectre est une véritable carte
d'identité chimique, lisible à des milliards de kilomètres.}
\end{blocEdition}

\editionPage{defi-cor-21-1}{D21 · Corrigé}{Combien de mesures faut-il ?}{\hyperref[defi-ex-21]{Énoncé p.~\pageref*{defi-ex-21}} · Démarche et vérifications}
\begin{blocEdition}{defi-cor-21-1}
\editionEtape{1. Des titrages simulés}
Chaque nombre est le résultat d'un titrage virtuel, entaché d'\res{\text{erreurs aléatoires}}. \code{random.seed(2026)} fixe le point de
départ du générateur : les tirages, donc l'affichage, sont \res{\text{les mêmes à chaque exécution}}, ce qui permet de vérifier les résultats.
\editionEtape{2. L'exécution}
\begin{console}
5 12.202 0.095 0.0425
20 12.194 0.091 0.0204
100 12.191 0.086 0.0086
500 12.203 0.08 0.0036
\end{console}
Chaque ligne donne $n$, la moyenne, $s$ et $u$, en mL. La moyenne reste proche de \qty{12,20}{mL} ; $s$ reste voisin de \qty{0,08}{mL}, car il décrit
la dispersion \res{\text{de chaque mesure}} ; seule $u$ diminue nettement : la moyenne de nombreuses mesures est plus précise.
\editionEtape{3. Cent fois plus de mesures}
$u=\dfrac{s}{\sqrt{n}}$ : si $n$ est multiplié par 100, $\sqrt{n}$ est multiplié par $\sqrt{100}=10$ et $u$ est \res{\text{divisée par 10}}. La simulation le confirme :
de 5 à 500 mesures, $u$ passe de \qty{0,0425}{mL} à \qty{0,0036}{mL}, aux fluctuations du hasard près.
\label{figedition-defi-cor-21-1-1}\schema[corrige]{defi-21-cor}
\end{blocEdition}

\editionPage{defi-cor-21-2}{D21 · Corrigé (suite)}{Combien de mesures faut-il ?}{\hyperref[defi-ex-21]{Énoncé p.~\pageref*{defi-ex-21}} · Démarche et vérifications}
\begin{blocEdition}{defi-cor-21-2}
\editionEtape{4. Le prix de la précision}
$n=\left(\dfrac{s}{u}\right)^2=\left(\dfrac{0{,}08}{0{,}001}\right)^2=\res{\num{6400}}$ titrages, soit $6400\times\qty{15}{min}=\qty{96000}{min}$ : environ \res{67~\text{jours}} de travail sans
interruption ! La répétition devient vite déraisonnable.
\editionEtape{5. Une erreur systématique}
Toutes les valeurs simulées sont décalées de \qty{0,10}{mL} : la moyenne tend vers \res{\qty{12,30}{mL}} et non vers \qty{12,20}{mL}, alors que $u$
continue de diminuer. Le résultat devient très précis\dots{} mais faux : répéter les mesures \res{\text{ne corrige pas}} une erreur systématique.
\editionEtape{6. Mieux que répéter}
Il vaut mieux \res{\text{réduire $s$}} : burette mieux réglée et lue à hauteur des yeux, ajout goutte à goutte près de l'équivalence,
indicateur au virage net ou suivi pH-métrique, concentrations choisies pour que $V_\text{éq}$ soit grand (l'incertitude relative diminue).
\editionRetenir{$u=\dfrac{s}{\sqrt{n}}$ : diviser l'incertitude-type par 10 coûte cent fois plus de mesures. Améliorer le protocole est souvent plus
efficace, et seule la vigilance permet de débusquer une erreur systématique.}
\end{blocEdition}
